\chapter{Příprava a analýza korpusu Toki Pona}
\label{chap:korpus}

V této kapitole představujeme metodiku sestavení textového korpusu, proces čištění a deduplikace a stratifikované rozdělení dat do trénovací, validační a testovací sady.

\section{Struktura a zdroje korpusu}
Pro zajištění žánrové a stylové rozmanitosti byl korpus vytvořen sloučením čtyř nezávislých otevřených zdrojů jazyka toki pona:
\begin{enumerate}
    \item \textbf{Tatoeba:} Databáze paralelních vět zaměřená na každodenní konverzační obraty a jednoduché větné struktury.
    \item \textbf{Wikipesija (Monolingual):} Encyklopedické články přeložené do toki pona, charakteristické složitějšími souvětími a popisným stylem.
    \item \textbf{Lipu Sewi:} Rozsáhlý překlad náboženských a klasických narativních textů s bohatým kompozicionálním vyjadřováním.
    \item \textbf{Poki Lapo:} Komunitní sborník tvůrčího psaní, povídek a poezie.
\end{enumerate}

\section{Čištění a deduplikace}
Surová data prošla automatizovaným předzpracováním:
\begin{itemize}
    \item převod veškerého textu na malá písmena (lowercase),
    \item kanonizace uvozovek a interpunkčních znamének,
    \item odstranění neplatných nelatinkových znaků mimo standardní abecedu toki pona,
    \item přesná deduplikace pomocí kryptografického hashe SHA256.
\end{itemize}

Původních 131\,434 vět bylo deduplikováno na výsledných **109\,885 unikátních vět**.

\section{Stratifikované rozdělení podle zdrojů}
Aby se zabránilo systematickému zkreslení validačních a testovacích odhadů (např. dominancí jednoho specifického žánru v testovacím souboru), bylo rozdělení provedeno stratifikovaně v poměru 80\,\% : 10\,\% : 10\,\% uvnitř každého dílčího zdroje:
\begin{table}[htbp]
    \centering
    \caption{Přehled rozdělení a velikostí korpusu toki pona}
    \label{tab:korpus_stats}
    \begin{tabular}{lrrr}
        \toprule
        Dílčí sada & Počet vět & Počet slov & Průměrná délka věty (znaky) \\
        \midrule
        Trénovací sada (80\,\%) & 87\,906 & 1\,415\,500 & 22{,}8 \\
        Validační sada (10\,\%) & 10\,987 & 177\,843 & 22{,}8 \\
        Testovací sada (10\,\%)  & 10\,992 & 175\,471 & 22{,}8 \\
        \midrule
        \textbf{Celkem} & \textbf{109\,885} & \textbf{1\,768\,814} & \textbf{22{,}8} \\
        \bottomrule
    \end{tabular}
\end{table}
